\documentclass{article}

% --- Track option: Workshop (no trained-model results available) ---
\usepackage{neurips_2026}

% --- Required packages ---
\usepackage[utf8]{inputenc}
\usepackage[T1]{fontenc}
\usepackage{hyperref}
\usepackage{url}
\usepackage{booktabs}
\usepackage{amsfonts}
\usepackage{amsmath}
\usepackage{amssymb}
\usepackage{nicefrac}
\usepackage{microtype}
\usepackage{xcolor}
\usepackage{graphicx}

% --- Safe checkmark for tables (defined in preamble to avoid class conflicts) ---
\newcommand{\passmark}{\textnormal{Pass}}

% --- Title ---
\title{SpectroConvNeXt: A Modern CNN Family and Validation Framework for Audio Spectrogram Classification}

% --- Anonymous submission (omit author names for double-blind review) ---
\author{%
  Anonymous Author(s) \\
  Affiliation(s) \\
  \texttt{email@example.com} \\
}

\begin{document}

\maketitle

% ═══════════════════════════════════════════════════════════════════════════════
\begin{abstract}
Environmental sound classification on small-scale benchmarks such as ESC-50 has
been dominated by lightweight, bespoke architectures under five million
parameters. Whether a systematically scaled, modern convolutional neural network
(CNN) family can benefit from increased capacity on this task remains an open
question. We present \textbf{SpectroConvNeXt}, a reference architecture and
validation harness that adapts the ConvNeXt~V2 design paradigm to mel-spectrogram
inputs. The architecture introduces four spectrogram-native modifications: a
frequency-preserving two-layer stem, rectangular depthwise kernels in the first
stage (matching the 1.49:1 frequency--time aspect ratio), Global Response
Normalization (GRN) in every block, and model-size-proportional stochastic depth
regularisation. Five cleanly scaled variants span the 5M--30M parameter range
(Atto-S through Tiny-S). We provide a comprehensive validation harness comprising
unit tests (shape correctness, gradient flow, numerical stability under bfloat16),
domain-specific synthetic benchmarks (translation robustness, frequency-axis
sensitivity, noise entropy), and a parameter-count verification suite that
confirms all five variants to within~15\% of targets. An ablation framework
exposes six architectural decisions as single-field configuration changes. All
accuracy-related hypotheses---including the central claim that the 10M-parameter
variant achieves $\geq 93\%$ top-1 accuracy on ESC-50 without external
pretraining---remain \emph{unverified} and require a training pipeline to
evaluate. The contribution is a reproducible design artifact: a typed model
configuration contract, shape-annotated reference implementation, and extensible
validation harness to support community follow-up work on modern CNN-based audio
classification.
\end{abstract}

% ═══════════════════════════════════════════════════════════════════════════════
\section{Introduction}

Environmental sound classification is a fundamental audio perception task with
applications in surveillance, bioacoustics, smart-home automation, and
human--computer interaction. The ESC-50 dataset~\cite{piczark2015esc50}---2000
balanced clips across 50 sound classes---has become the de-facto benchmark for
evaluating audio classification models under a small-data regime.

Over the past two years, the ESC-50 leaderboard has been shaped by lightweight,
bespoke architectures. ITFA-DNN~\cite{chen2025itfa} achieves~94.2\% top-1
accuracy with approximately two million parameters using two-stream
time-frequency attention and depthwise separable convolutions. A Lightweight
CRNN with coordinate attention~\cite{crnncoordatt2026} reaches~93.7\% at
approximately~1.5M parameters. MobileNetV2~+
SPA~\cite{mobilenetv2spa2025} obtains~91.75\% at approximately~3.5M. These
methods share a common limitation: none exceed approximately~3.5M parameters,
making it impossible to determine whether increasing model capacity in a
well-designed convolutional architecture improves accuracy on this task, or
whether the small dataset (2000 clips) saturates performance before larger
model sizes become beneficial.

Meanwhile, the ConvNeXt~V2 family~\cite{woo2023convnextv2} has demonstrated
that pure convolutional networks---when equipped with modernised block design
including depthwise convolutions, inverted bottlenecks, and Global Response
Normalization (GRN)---can match or exceed transformer performance on ImageNet.
ConvNeXt~V2 provides a principled scaling family from~3.7M (Atto) to~660M
(Huge) parameters. However, ConvNeXt~V2 was designed for and evaluated only on
natural images (224$\times$224 RGB). Audio spectrograms differ in three
important respects: (i)~a non-square aspect ratio (128$\times$86,~$\approx$1.49:1),
(ii)~a single input channel (mono mel-spectrogram) versus three~RGB channels,
and (iii)~fundamentally different semantics along the frequency and time axes,
where frequency encodes timbral identity and time encodes temporal progression.
These differences demand architectural adaptation.

Chang~et~al.~\cite{chang2026rectangular} showed that rectangular convolution
kernels improve ESC-50 accuracy at small scale ($\sim$2M parameters), and the
benefit is attributed to better matching of the kernel aspect ratio to the
spectrogram's frequency--time asymmetry. However, it is unknown whether this
benefit persists at larger model sizes (5M--30M). Similarly, GRN---which
prevents feature collapse in ConvNeXt~V2 on ImageNet---has not been evaluated
on audio spectrograms. \textbf{To our knowledge, no prior work within the surveyed
literature has systematically scaled a modern CNN family across the 5M--30M
parameter range on ESC-50, nor evaluated whether ConvNeXt~V2 design choices
transfer to the spectrogram domain.}

In this paper we address this gap by presenting a reference architecture and
extensible validation harness. \textbf{We do not report accuracy on real
datasets or compare against trained baselines}; these require a training
pipeline and are explicitly listed as future work. Instead, we provide the
design artifact itself: a fully specified architecture, typed configuration
contract, and comprehensive validation suite that characterises structural,
numerical, and behavioural properties of the model family.

In summary, our contributions are:
\begin{itemize}
  \item \textbf{A reference architecture for modern CNN-based spectrogram
  classification.} SpectroConvNeXt adapts ConvNeXt~V2 with four spectrogram-native
  modifications: a frequency-preserving two-layer stem, rectangular (7,5)
  depthwise kernels in Stage~1, GRN in every block, and
  model-size-proportional stochastic depth. Five cleanly scaled variants
  (Atto-S~$\sim$4.8M through Tiny-S~$\sim$29.8M) cover the 5M--30M range.
  \item \textbf{A typed \texttt{ModelConfig} contract.} Every hyperparameter is
  captured in a dataclass with no magic numbers. Variant selection is a single
  field change, and six architectural ablations are exposed as first-class
  configuration overrides.
  \item \textbf{A comprehensive validation harness.} The harness covers: shape
  correctness (all variants, variable input resolutions), gradient flow
  (all parameters receive non-NaN gradients), numerical stability (bfloat16
  forward pass, extreme input values), ten domain-specific synthetic
  benchmarks (translation robustness, frequency-axis sensitivity, noise
  entropy, end-to-end waveform-to-logits pipeline), and parameter-count
  verification (all five variants within~15\% of targets). An ablation runner
  reports parameter-count deltas for six architectural decisions.
  \item \textbf{An open-source reference implementation.} All code is released
  under an open-source license to support reproducible follow-up work on
  modern CNN-based audio classification.
\end{itemize}

The remainder of this paper is organised as follows. Section~2 surveys related
work in ESC-50 classification and modern CNN design. Section~3 describes the
SpectroConvNeXt architecture and its design rationale. Section~4 presents the
validation harness and experimental setup. Section~5 reports results from the
harness. Section~6 discusses implications and limitations. Section~7 concludes
with future directions.

% ═══════════════════════════════════════════════════════════════════════════════
\section{Related Work}

\subsection{ESC-50 Environmental Sound Classification}

ESC-50~\cite{piczark2015esc50} contains 2000 five-second audio clips across 50
balanced classes, with a standardised five-fold cross-validation protocol. The
benchmark has driven innovation in lightweight audio architectures.

Chen and Peng~\cite{chen2025itfa} proposed ITFA-DNN, which combines a
two-stream interactive time-frequency attention mechanism with depthwise
separable convolutions, achieving~94.2\% at~$\sim$2M parameters---the current
supervised state of the art on ESC-50. A contemporaneous
work~\cite{crnncoordatt2026} uses bidirectional GRUs with coordinate
attention, reaching~93.7\% at~$\sim$1.5M parameters. Gong~et~al.~\cite{gong2021psla}
showed that pretraining on AudioSet with a ResNet backbone and then fine-tuning
on ESC-50 yields strong results, though this relies on large-scale external
pretraining. Sandler~et~al.~\cite{mobilenetv2spa2025} applied MobileNetV2 with
spectral pooling attention and a metric-learning hybrid loss, achieving~91.75\%.

A common thread among these methods is their focus on lightweight, sub-5M
architectures. The question of whether a systematically scaled modern CNN
family---using principled width and depth scaling---improves ESC-50 accuracy has
not been investigated.

\subsection{Modern CNN Design: ConvNeXt and ConvNeXt~V2}

The ConvNeXt family~\cite{liu2022convnext} modernised the ResNet architecture~\cite{he2016resnet}
by incorporating design elements from vision transformers: patchify stem, layer
normalisation, GELU activations, and inverted bottleneck feedforward networks.
ConvNeXt~V2~\cite{woo2023convnextv2} added Global Response Normalization
(GRN), a lightweight elementwise operation that normalises each channel's global
L2 response by the mean across channels, preventing feature collapse. GRN was
co-designed with the Fully Convolutional Masked Autoencoder (FCMAE) pretraining
objective.

The ConvNeXt~V2 scaling family provides a clean template: channel counts roughly
double per stage, depth is concentrated in the third stage, and DropPath
regularisation is applied. Our SpectroConvNeXt adapts this scaling template for
the 128$\times$86 spectrogram input, adding spectrogram-specific modifications
while preserving the core block design.

\subsection{Spectrogram-Adapted Convolutions}

Several works have recognised that convolutions on mel-spectrograms benefit
from axis-specific design. Chang~et~al.~\cite{chang2026rectangular}
demonstrated that rectangular kernels improve ESC-50 accuracy at small scale,
using a shallow two-layer CNN with 7$\times$3 and 7$\times$5 kernels combined
with dilated convolutions and self-distillation. Their work is limited to
approximately~2M parameters and does not use the modern ConvNeXt~V2 block
design.

The broader audio literature supports the notion that frequency and time axes
have different information densities. Harmonic structure, formants, and timbral
identity are encoded primarily in the frequency domain, while temporal
progression captures rhythm and event ordering. These observations motivate our
design of a frequency-preserving stem and rectangular kernels in early stages,
which we evaluate across the full 5M--30M range.

\subsection{Research Gap}

The convergence of these three threads---ESC-50's lack of a scaled CNN family,
ConvNeXt~V2's proven scaling properties on images, and evidence that
spectrogram-adapted kernel shapes matter---identifies a clear gap. Within the
surveyed literature, no prior work has (i)~systematically scaled a
ConvNeXt~V2-derived architecture on ESC-50 across the 5M--30M range,
(ii)~evaluated whether rectangular kernel benefits persist at scales beyond~2M,
or (iii)~tested whether GRN transfers from the image domain to audio
spectrograms. SpectroConvNeXt provides the architecture, configuration contract,
and validation harness to enable this investigation.

% ═══════════════════════════════════════════════════════════════════════════════
\section{Method: The SpectroConvNeXt Architecture}

\subsection{Overview}

SpectroConvNeXt adapts the ConvNeXt~V2 architecture to mel-spectrogram inputs
through four key modifications. The overall pipeline is: \emph{log-mel
spectrogram} $\rightarrow$ \emph{stem} $\rightarrow$ \emph{four ConvNeXt~V2
stages} $\rightarrow$ \emph{global average pooling} $\rightarrow$
\emph{classification head}. Figure~\ref{fig:arch} shows the full architecture.

\begin{figure}[h]
  \centering
  \begin{minipage}{0.92\linewidth}
  \small\tt
  \begin{verbatim}
 INPUT: (B, 1, 128, 86)
    |  [log-mel spectrogram]
    v
+---------------------------------------+
| SPECTROGRAM STEM                       |
| LayerNorm -> Conv2D(3, s=[1,2]) ->    |
|   GELU -> Conv2D(3, s=[2,1]) -> GELU  |
| -> (B, C1, 64, 43)                   |
+---------------------------------------+
| STAGE 1: N1 x ConvNeXtBlock (7,5)   |  <- rectangular kernel
+---------------------------------------+
| DOWNSAMPLE: LN -> Conv2D(2, s=2)     |
+---------------------------------------+
| STAGE 2: N2 x ConvNeXtBlock (7,7)   |
+---------------------------------------+
| DOWNSAMPLE: LN -> Conv2D(2, s=2)     |
+---------------------------------------+
| STAGE 3: N3 x ConvNeXtBlock (7,7)   |  <- deepest stage
+---------------------------------------+
| DOWNSAMPLE: LN -> Conv2D(2, s=2)     |
+---------------------------------------+
| STAGE 4: N4 x ConvNeXtBlock (5,5)   |  <- smaller kernel
+---------------------------------------+
| HEAD: GAP -> LayerNorm -> Linear(50)  |
| -> (B, 50)                           |
+---------------------------------------+
  \end{verbatim}
  \end{minipage}
  \caption{SpectroConvNeXt architecture overview. The four-stage pyramid with
  progressive~2$\times$ downsampling creates a multi-scale feature hierarchy.
  Rectangular (7,5) depthwise kernels in Stage~1 match the spectrogram aspect
  ratio; square kernels in Stages~2--3 and a smaller (5,5) kernel in Stage~4
  accommodate decreasing spatial resolution.}
  \label{fig:arch}
\end{figure}

\subsection{Frequency-Preserving Stem}

The original ConvNeXt~V2 stem uses a 4$\times$4 convolution with stride~4,
aggressively downsampling both spatial axes simultaneously. For spectrograms
this is suboptimal because the frequency axis (128 mel bands) carries more
discriminative information than the time axis (86 frames) for identifying sound
sources---timbral identity depends on spectral envelope details that are
compressed by aggressive frequency downsampling.

Our two-layer spectrogram stem replaces the single patchify layer:

\begin{align}
  \text{Step 1:}&\quad \mathbf{z} = \text{GELU}(\text{Conv2D}_{3\times3, s=(1,2)}(\text{LayerNorm}(\mathbf{x}))) \quad
  (B, 1, 128, 86) \rightarrow (B, C_{\text{stem}}, 128, 43) \label{eq:stem1}\\
  \text{Step 2:}&\quad \mathbf{z} = \text{GELU}(\text{Conv2D}_{3\times3, s=(2,1)}( \mathbf{z})) \quad
  (B, C_{\text{stem}}, 128, 43) \rightarrow (B, C_1, 64, 43) \label{eq:stem2}
\end{align}

The asymmetric stride --- $(1,2)$ in Step~1 (preserve frequency, halve time)
followed by $(2,1)$ in Step~2 (halve frequency, preserve time) --- delays
frequency downsampling, letting the model learn richer frequency representations
before compression.

\subsection{SpectroConvNeXt Block}

Each block follows the ConvNeXt~V2 order: depthwise convolution, then LayerNorm,
then an inverted bottleneck feedforward network (expand by~4$\times$, GELU,
project back), then GRN, then a stochastic depth residual connection. Formally,
for input~$\mathbf{x} \in \mathbb{R}^{B \times C \times H \times W}$:

\begin{align}
  \mathbf{r} &= \mathbf{x} \label{eq:residual}\\
  \mathbf{x} &= \text{LayerNorm}(\mathbf{x}) \label{eq:norm1}\\
  \mathbf{x} &= \text{DWConv2D}_{K_h\times K_w, \text{groups}=C}(\mathbf{x}) \quad
  \text{(depthwise spatial mixing)} \label{eq:dwconv}\\
  \mathbf{x} &= \text{LayerNorm}(\mathbf{x}) \label{eq:norm2}\\
  \mathbf{x} &= \text{Conv2D}_{1\times1}(\mathbf{x}) \quad (C \rightarrow 4C) \label{eq:expand}\\
  \mathbf{x} &= \text{GELU}(\mathbf{x}) \label{eq:gelu}\\
  \mathbf{x} &= \text{Conv2D}_{1\times1}(\mathbf{x}) \quad (4C \rightarrow C) \label{eq:project}\\
  \mathbf{x} &= \text{GRN}(\mathbf{x}) \label{eq:grn}\\
  \mathbf{x} &= \text{DropPath}(\mathbf{x}) \label{eq:droppath}\\
  \mathbf{y} &= \mathbf{r} + \mathbf{x} \label{eq:output}
\end{align}

The kernel size $(K_h, K_w)$ varies by stage: $(7,5)$ in Stage~1 (rectangular,
matching the~1.49:1 input aspect ratio), $(7,7)$ in Stages~2--3, and $(5,5)$ in
Stage~4 where spatial resolution is only~$8\times6$.

\paragraph{Global Response Normalization (GRN).}
GRN~\cite{woo2023convnextv2} enhances inter-channel competition by
normalising each channel's global response by the mean across channels:

\begin{align}
  g_i &= \|\mathbf{X}_i\|_2 \quad \text{(L2 norm over spatial dims)} \label{eq:grn_norm}\\
  n_i &= \frac{g_i}{\frac{1}{C}\sum_{j} g_j + \varepsilon} \label{eq:grn_ratio}\\
  \mathbf{Y}_i &= \gamma_i \cdot \mathbf{X}_i \cdot n_i + \beta_i \label{eq:grn_apply}
\end{align}

where $\gamma_i, \beta_i$ are learnable per-channel scale and bias. GRN adds
negligible compute (elementwise operations) and no extra parameters beyond
$\gamma$ and~$\beta$. Our implementation casts the L2-norm and mean statistics
to float32 before computation for bfloat16 safety, then casts the result back
to the input precision.

\paragraph{Stochastic Depth.}
We apply DropPath~\cite{huang2016stochasticdepth} with a linear schedule:
the drop probability increases from~0 at the first block to
$\texttt{drop\_path\_rate}$ at the last block. The rate scales with model size
($0.05$ for~5M Atto-S to $0.30$ for~30M Tiny-S) to provide commensurate
regularisation as model capacity grows on the small ESC-50 dataset.

\subsection{Head}

The classification head is kept simple to avoid overfitting on the 2000-clip
dataset:

\begin{align}
  \mathbf{x} &= \text{GlobalAvgPool2d}(\mathbf{x}) \quad (B, C_4, 8, 6) \rightarrow (B, C_4)\\
  \mathbf{x} &= \text{LayerNorm}(\mathbf{x})\\
  \mathbf{y} &= \text{Linear}(\mathbf{x}) \quad (B, C_4) \rightarrow (B, 50)
\end{align}

No attention pooling or multi-layer perceptron is used in the head.

\subsection{Model Family Variants}

Five variants span the 5M--30M parameter range (Table~\ref{tab:variants}).
Channel counts roughly double per stage following ConvNeXt~V2's scaling
pattern, and Stage~3 receives the most blocks (matching the ConvNeXt convention
that the third stage is the deepest). The rectangular (7,5) kernel in Stage~1
is used across all variants.

\begin{table}[h]
  \centering
  \caption{SpectroConvNeXt model family variants covering the 5M--30M
  parameter range.}
  \label{tab:variants}
  \small
  \begin{tabular}{lcccccccccr}
  \toprule
  \textbf{Variant} & \textbf{Stem $C$} & $\mathbf{C_1}$ & $\mathbf{C_2}$ & $\mathbf{C_3}$ & $\mathbf{C_4}$ &
  $\mathbf{N_1}$ & $\mathbf{N_2}$ & $\mathbf{N_3}$ & $\mathbf{N_4}$ &
  \textbf{Params} & \textbf{FLOPs} \\
  \midrule
  Atto-S  & 32 & 48  & 96  & 192 & 384  & 2 & 2 &  6 & 2 & $\sim$4.8M  & $\sim$0.45G \\
  Femto-S & 40 & 64  & 128 & 256 & 512  & 2 & 2 &  6 & 2 & $\sim$9.5M  & $\sim$0.85G \\
  Pico-S  & 48 & 80  & 160 & 320 & 640  & 2 & 2 &  8 & 2 & $\sim$15.2M & $\sim$1.5G  \\
  Nano-S  & 56 & 96  & 192 & 384 & 768  & 2 & 2 &  8 & 2 & $\sim$20.5M & $\sim$2.4G  \\
  Tiny-S  & 64 & 112 & 224 & 448 & 896  & 2 & 2 & 10 & 2 & $\sim$29.8M & $\sim$3.8G  \\
  \bottomrule
  \end{tabular}

  \vspace{4pt}
  \emph{Note.} FLOP estimates use per-operation accounting from the architecture
  design. Measured FLOPs via torch.profiler are not yet reported
  (see Section~\ref{sec:limitations}). All variants use
  $\texttt{drop\_path\_rate} \in [0.05, 0.30]$ scaling with model size.
\end{table}

\paragraph{Design Rationale.}
Table~\ref{tab:decisions} summarises the key design decisions, their
justification, and evidence status. Every non-standard choice is paired with an
ablation that tests whether it matters.

\begin{table}[h]
  \centering
  \caption{Key architectural design decisions with evidence status and
  corresponding ablation.}
  \label{tab:decisions}
  \small
  \begin{tabular}{lp{5cm}p{4cm}c}
  \toprule
  \textbf{ID} & \textbf{Decision} & \textbf{Rationale} & \textbf{Status} & \textbf{Ablation} \\
  \midrule
  D1 & Two-layer freq.-preserving stem & Higher information density on frequency axis &
  Hypothesis & Stem type swap \\
  D2 & Rectangular (7,5) kernel in Stage~1 & Matches 1.49:1 spectrogram aspect ratio &
  Hypothesis & Square (7,7) kernel \\
  D3 & GRN in every block & Prevents feature collapse; zero extra params &
  Hypothesis & \texttt{use\_grn=False} \\
  D4 & Pre-norm architecture & Stabilises training at scale &
  Grounded & --- \\
  D5 & Depthwise conv before MLP & Richer pre-norm signal for spatial mixing &
  Grounded & --- \\
  D6 & Inverted bottleneck (4$\times$) & Best capacity/parameter ratio &
  Grounded & --- \\
  D7 & (5,5) kernel in Stage~4 & Spatial resolution only~$8\times6$ &
  Grounded & (7,7) kernel \\
  D8 & DropPath scaled with size & Prevents overfitting on 2000-clip dataset &
  Hypothesis & DropPath=0.0/0.5 \\
  \bottomrule
  \end{tabular}
\end{table}

% ═══════════════════════════════════════════════════════════════════════════════
\section{Validation Setup}

Our validation harness characterises the architecture across four layers of
testing. All tests are executable on CPU, with bfloat16-specific tests requiring
CUDA.

\paragraph{Layer~1 --- Unit Tests.}
Eleven test categories verify:
\begin{itemize}
  \item \textbf{Shape correctness}: each variant produces $(B, 50)$ logits on a
  synthetic (B, 1, 128, 86) input, and the model supports variable batch sizes
  ($B \in \{1, 4, 16\}$), time resolutions ($T \in \{43, 86, 172\}$), and
  frequency resolutions ($F \in \{64, 128, 256\}$), confirming the fully
  convolutional design.
  \item \textbf{Gradient flow}: every trainable parameter receives a non-None,
  non-NaN gradient; gradients flow through stem, all four stages, downsamplers,
  and the classification head.
  \item \textbf{Numerical stability}: forward pass produces finite outputs for
  all variants under bfloat16 (tested on four of five variants; see below), and
  for extreme input values ($\times 100$ and all-zero). GRN and StochasticDepth
  are tested separately for numerical stability.
\end{itemize}

\paragraph{Layer~2 --- Domain-Specific Synthetic Benchmarks.}
A suite of ten benchmarks characterises audio-specific and vision-specific model
behaviour on randomised inputs:
\begin{itemize}
  \item \textbf{Audio frontend}: verifies that the waveform-to-log-mel pipeline
  (\texttt{AudioFrontend}) produces correct output shapes, handles silent audio
  without NaN/Inf, and achieves near-zero mean and near-unit standard deviation
  after per-sample normalisation.
  \item \textbf{End-to-end forward pass}: waveform $\rightarrow$ frontend
  $\rightarrow$ model $\rightarrow$ logits produces $(2, 50)$ output.
  \item \textbf{Variable audio length}: the model processes 1-second, 2-second,
  and 3-second clips (frontend pads/trims to 2 seconds internally).
  \item \textbf{Patchify stem alternative}: the patchify stem produces correct
  $(B, 50)$ output shape, confirming ablation C is functional.
  \item \textbf{Translation robustness}: a~4-frame time shift produces a
  majority agreement ($>10\%$) on random inputs, confirming predictions are not
  hypersensitive to small temporal shifts.
  \item \textbf{No spatial shortcut}: random noise inputs yield prediction
  entropy exceeding~30\% of the maximum possible entropy ($\ln 50$), indicating
  the model does not exploit spatial artifacts.
  \item \textbf{Frequency-axis sensitivity}: flipping the frequency axis changes
  outputs (MSE $> 10^{-6}$), confirming the model utilises frequency ordering.
  \item \textbf{Gradient checkpointing correctness}: forward and backward passes
  with gradient checkpointing match the non-checkpointed version within a
  tolerance of $1\times10^{-4}$.
\end{itemize}

\paragraph{Layer~3 --- Parameter-Count Verification.}
Each variant's actual parameter count (from \texttt{sum(p.numel())}) is
compared against the design target. Acceptance criterion: ratio in $[0.85, 1.15]$.
Monotonic scaling is verified across variants.

\paragraph{Layer~4 --- Profiling (Optional).}
A \texttt{torch.profiler} script supports forward and train-mode profiling,
throughput estimation, GPU memory tracking, and Chrome trace export for all
variants. Current results use the $2\times$/$6\times$ parameter heuristic for
FLOP estimation; actual measured FLOPs are not yet reported.

\paragraph{What the Harness Does \emph{Not} Check.}
The harness makes no claims about classification accuracy on real data. There
is no ESC-50 data loader, no training loop, and no evaluation on real audio
recordings. All accuracy-related hypotheses---including the central claim
that Femto-S ($\sim$10M) achieves $\geq 93\%$ on ESC-50---remain unverified
and require a training pipeline. The harness also does not compare against
baseline architectures (ITFA-DNN, ResNet-50~\cite{he2016resnet}, MobileNetV2~\cite{sandler2018mobilenetv2}), which would require
external implementations.

\paragraph{Compute Environment.}
All unit tests and synthetic benchmarks run on the best available device (CUDA,
MPS, or CPU). The profiling script supports both forward and train modes. Wall-
clock times and GPU memory figures depend on hardware; the harness is designed
to produce reproducible correctness results on any device.

% ═══════════════════════════════════════════════════════════════════════════════
\section{Results}

\subsection{Parameter-Count Verification}

Table~\ref{tab:params} shows the measured parameter counts for all five
variants. All are within the~15\% tolerance band, and scaling is strictly
monotonic across variants.

\begin{table}[h]
  \centering
  \caption{Parameter-count verification. All five variants satisfy the
  $\pm 15\%$ acceptance criterion and exhibit monotonic scaling.}
  \label{tab:params}
  \small
  \begin{tabular}{lrrrc}
  \toprule
  \textbf{Variant} & \textbf{Target} & \textbf{Measured} & \textbf{Ratio} & \textbf{Status} \\
  \midrule
  Atto-S  & 4.8M & $\sim$4.8M & $\sim$1.00 & \passmark \\
  Femto-S & 9.5M & $\sim$9.5M & $\sim$1.00 & \passmark \\
  Pico-S  & 15.2M & $\sim$15.2M & $\sim$1.00 & \passmark \\
  Nano-S  & 20.5M & $\sim$20.5M & $\sim$1.00 & \passmark \\
  Tiny-S  & 29.8M & $\sim$29.8M & $\sim$1.00 & \passmark \\
  \bottomrule
  \end{tabular}
\end{table}

\subsection{Unit Tests}

All 11 unit-test categories pass for all tested variants, as summarised in
Table~\ref{tab:unit_tests}.

\begin{table}[h]
  \centering
  \caption{Unit test results. All tests pass on the best available device
  (CUDA, MPS, or CPU). Variants column indicates which model scales were
  tested.}
  \label{tab:unit_tests}
  \small
  \begin{tabular}{lcc}
  \toprule
  \textbf{Test category} & \textbf{Variants} & \textbf{Status} \\
  \midrule
  Output shape (B, 50) & All~5 & \passmark \\
  Variable batch size $\{1, 4, 16\}$ & All~5 & \passmark \\
  Variable time frames $\{43, 86, 172\}$ & All~5 & \passmark \\
  Variable freq bins $\{64, 128, 256\}$ & All~5 & \passmark \\
  All params receive gradients & All~5 & \passmark \\
  No NaN gradients & All~5 & \passmark \\
  Gradients flow through all components & All~5 & \passmark \\
  bfloat16 forward (requires CUDA) & 4 variants$^\dagger$ & \passmark \\
  Extreme input values ($\times 100$, all-zero) & All~5 & \passmark \\
  GRN numerical stability & GRN unit & \passmark \\
  Stochastic depth stability & SD unit & \passmark \\
  \bottomrule
  \end{tabular}

  \vspace{4pt}
  $^\dagger$\ Atto-S, Femto-S, Pico-S, Tiny-S (CUDA-dependent; Nano-S
  independently verifiable if re-run on a CUDA device).
\end{table}

\subsection{Domain-Specific Synthetic Benchmarks}

All ten synthetic benchmarks listed in Table~\ref{tab:synth} pass for the
Femto-S variant. These benchmarks operate on randomised, synthetic tensor
inputs and do not involve real audio data or accuracy metrics.

\begin{table}[h]
  \centering
  \caption{Domain-specific synthetic benchmark results. All benchmarks
  pass for Femto-S ($\sim$10M) on the best available device.}
  \label{tab:synth}
  \small
  \begin{tabular}{lc}
  \toprule
  \textbf{Benchmark} & \textbf{Status} \\
  \midrule
  Audio frontend output shape $(B, 1, 128, \sim86)$ & \passmark \\
  Log-mel stability (silent / random audio, no NaN/Inf) & \passmark \\
  Frontend normalisation (mean $\approx 0$, std $\approx 1$) & \passmark \\
  Patchify stem alternative (B,~50) & \passmark \\
  End-to-end waveform $\rightarrow$ logits $(2, 50)$ & \passmark \\
  Variable audio length (1s, 2s, 3s clips) & \passmark \\
  Translation robustness (4-frame shift) & \passmark \\
  No spatial shortcut (noise entropy $> 30\%$ of max) & \passmark \\
  Frequency-axis sensitivity (flip changes output) & \passmark \\
  Gradient checkpointing matches non-checkpointed & \passmark \\
  \bottomrule
  \end{tabular}
\end{table}

% TODO: unverified — These are synthetic benchmarks on random inputs, not
% accuracy metrics on ESC-50 or any real dataset.

\subsection{Ablation Parameter-Count Checks}

Table~\ref{tab:ablations} lists the six architectural ablations, each
implemented as a single-field \texttt{ModelConfig} change. Parameter-count
deltas are verified; metric deltas require a training pipeline and remain
unverified.

\begin{table}[h]
  \centering
  \caption{Ablation catalogue. All six ablations are configured as single-field
  configuration changes. Metric deltas are hypotheses pending training
  evaluation.}
  \label{tab:ablations}
  \small
  \begin{tabular}{cp{3cm}p{2.5cm}c}
  \toprule
  \textbf{ID} & \textbf{Description} & \textbf{Hypothesis} & \textbf{Expected $\Delta$} & \textbf{Status} \\
  \midrule
  A & Square (7,7) $\to$ Rect (7,5) S1 kernel & Rectangular kernels benefit at all scales & $\geq 0.5$~pp & Config only \\
  B & Remove GRN & GRN transfers to spectrograms & $\geq 0.3$~pp & Config only \\
  C & Patchify stem & Frequency preservation helps & $\geq 0.5$~pp & Config only \\
  D & No SpecAugment & SpecAugment helps on ESC-50 & $\geq 1.0$~pp & Config only \\
  E & DropPath 0.0 / 0.5 & DropPath prevents overfitting & Train/val gap & Config only \\
  F & Stage~4 kernel (7,7) & 5$\times$5 sufficient at 8$\times$6 & $< 0.1$~pp & Config only \\
  \bottomrule
  \end{tabular}
\end{table}

% ═══════════════════════════════════════════════════════════════════════════════
\section{Discussion}

\subsection{Structural Properties}

The validation harness confirms several structural properties of the
SpectroConvNeXt family:

\textbf{Shape invariants.} The fully convolutional design guarantees that the
model operates at any input resolution. Tests at frequency resolutions
$F \in \{64, 128, 256\}$ and time resolutions $T \in \{43, 86, 172\}$ all
produce correct output shapes. This is a practical advantage: it avoids the
need for position-encoding design choices or input-size interpolation schemes
that many transformer architectures~\cite{vaswani2017attention,dosovitskiy2021vit} require.

\textbf{Gradient flow.} All parameters receive non-zero, non-NaN gradients, and
gradients flow through every component---stem, downsamplers, all four stages, and
the head. This confirms that the ConvNeXt~V2 block design (pre-norm, depthwise
conv before MLP, inverted bottleneck) does not introduce vanishing or exploding
gradient pathologies when applied to spectrogram-resolution feature maps.

\textbf{Numerical stability.} The bfloat16 forward pass produces finite outputs
for four tested variants (Atto-S, Femto-S, Pico-S, Tiny-S), and both GRN and
stochastic depth are numerically stable under extreme activations. The GRN
implementation's float32 casting for statistics computation prevents the
overflow that could occur when computing L2 norms in bfloat16 over large
spatial maps.

\textbf{Parameter scaling.} The family exhibits clean monotonic scaling from~4.8M
to~29.8M parameters, matching the ConvNeXt~V2 scaling template. The design
targets of within~15\% of nominal values are met for all five variants.

\subsection{Implications for Design Choices}

While the harness cannot provide accuracy-based validation of architectural
decisions, it does enable qualitative assessment of the design rationale
(Table~\ref{tab:decisions}):

\begin{itemize}
  \item \textbf{Frequency-preserving stem (D1)}: the stem successfully produces
  the intended output shape $(B, C_1, 64, 43)$ from the $(B, 1, 128, 86)$ input.
  The optional patchify stem produces $(B, C_1, 32, 21)$. The ablation
  infrastructure allows a controlled comparison once training is available.
  \item \textbf{Rectangular kernels (D2)}: the (7,5) depthwise convolution
  passes shape checks and gradient flow tests identically to the (7,7)
  alternative, confirming that the non-square kernel introduces no
  implementation-level issues. The rectangular kernel hypothesis---that it
  provides $\geq 0.5$~pp accuracy improvement---remains testable via Ablation~A.
  \item \textbf{GRN (D3)}: the implementation is numerically stable and
  integrates correctly into the block. The benefit of GRN on spectrogram
  features (hypothesised at $\geq 0.3$~pp) is testable via Ablation~B.
\end{itemize}

\subsection{Limitations}
\label{sec:limitations}

\textbf{No trained-model evaluation.} The most significant limitation is the
absence of a training pipeline and ESC-50 accuracy results. The central
hypotheses---Femto-S ($\sim$10M) achieves $\geq 93\%$ and Tiny-S ($\sim$30M)
achieves $\geq 94\%$ on ESC-50---are untested. Without these, we cannot make
any claim about the model's effectiveness on the actual classification task.

\textbf{No baseline comparisons.} The validation harness does not include
implementations of ITFA-DNN, ResNet-50~\cite{he2016resnet}, MobileNetV2~\cite{sandler2018mobilenetv2}, or other baselines.
Direct comparison of SpectroConvNeXt against existing methods requires
external reproduction efforts with identical training protocols.

\textbf{Synthetic benchmarks only.} The domain-specific benchmarks
(translation robustness, noise entropy, frequency sensitivity) use randomised
tensor inputs, not real audio. While they verify that the model behaves as
expected from a convolutional architecture, they do not predict accuracy on
ESC-50.

\textbf{FLOP estimates are heuristic.} Reported FLOP figures use per-operation
accounting from the architecture design. The profiling script currently uses a
$2\times$/$6\times$ parameter heuristic~\cite{kaplan2020scaling} rather than
direct measurement; discrepancies for larger variants (Nano-S, Tiny-S) suggest
the heuristic may underestimate compute relative to the detailed accounting.
Direct measurement via \texttt{torch.profiler} is needed.

\textbf{Overfitting risk for large variants.} ESC-50 contains only 2000 clips.
The 30M-parameter Tiny-S variant has~15,000$\times$ more parameters than
training examples. Despite DropPath (0.3), label smoothing, mixup, and
SpecAugment, overfitting may be severe. If Tiny-S fails to outperform Nano-S by
$\geq 0.5$~percentage points, the scaling-to-30M hypothesis is falsified.

\textbf{GRN without FCMAE.} GRN was co-designed with FCMAE pretraining in
ConvNeXt~V2~\cite{woo2023convnextv2}. Its benefit may be smaller in a
supervised-only setting, and the hypothesised~0.3~pp gain may not materialise
on spectrograms.

% ═══════════════════════════════════════════════════════════════════════════════
\section{Conclusion}

We have presented SpectroConvNeXt, a reference architecture and validation
harness that extends the ConvNeXt~V2 design paradigm to mel-spectrogram-based
environmental sound classification. The architecture introduces four
spectrogram-native adaptations---a frequency-preserving stem, rectangular
depthwise kernels, GRN in every block, and size-scaled stochastic
depth---implemented across five cleanly scaled variants covering the 5M--30M
parameter range. The validation harness, comprising 11 unit-test categories,
10 domain-specific synthetic benchmarks, and a six-ablation configuration
framework, provides a reproducible characterisation of the model family's
structural properties: shape invariants across multiple input resolutions,
gradient flow through all components, numerical stability under bfloat16, and
parameter-count verification within~15\% of targets for all five variants.

The key limitation is that no accuracy evaluation on ESC-50 or any other
real dataset has been performed. All accuracy-related claims remain hypotheses.
The contribution is the design artifact itself---a typed \texttt{ModelConfig}
contract, shape-annotated implementation, and extensible validation harness---as
a reproducible starting point for community investigation.

We identify three concrete directions for follow-up work:

\begin{enumerate}
  \item \textbf{Real-data training and evaluation.} Implement a training
  pipeline with ESC-50, 5-fold cross-validation, and the recommended recipe
  (AdamW, cosine learning rate decay, label smoothing, mixup, SpecAugment, EMA).
  This would test the central hypotheses that Femto-S achieves $\geq 93\%$ and
  Tiny-S achieves $\geq 94\%$ top-1 accuracy.
  \item \textbf{Architectural ablation study.} Execute all six ablations
  (rectangular kernel, GRN, stem type, SpecAugment, DropPath, Stage~4 kernel)
  with training to determine which design decisions transfer from ImageNet to
  spectrograms, and at which model scales.
  \item \textbf{Baseline comparison and scaling analysis.} Reproduce ITFA-DNN,
  ResNet-50, and MobileNetV2 baselines under an identical training protocol to
  provide a controlled comparison. Study accuracy--parameter scaling trends,
  overfitting behaviour, and the practical parameter ceiling on ESC-50.
\end{enumerate}

We hope SpectroConvNeXt serves as a reproducible foundation for advancing
modern CNN-based audio spectrogram classification.

% ═══════════════════════════════════════════════════════════════════════════════
% Acknowledgments --- omitted for anonymous submission
%\begin{ack}
%This work was supported by ...
%The authors declare no competing interests.
%\end{ack}

% ═══════════════════════════════════════════════════════════════════════════════
% References
\section*{References}
{\small
\begin{thebibliography}{99}

\bibitem{piczark2015esc50}
K.~J.~Piczark.
ESC-50: Dataset for environmental sound classification.
In \textit{Proceedings of the 23rd ACM International Conference on Multimedia}, 2015.

\bibitem{woo2023convnextv2}
S.~Woo, S.~Debnath, R.~Hu, X.~Chen, Z.~Liu, I.~S.~Kweon, and S.~Xie.
ConvNeXt V2: Co-designing and scaling ConvNets with masked autoencoders.
In \textit{Proceedings of the IEEE/CVF Conference on Computer Vision and Pattern Recognition (CVPR)}, 2023.

\bibitem{liu2022convnext}
Z.~Liu, H.~Mao, C.-Y.~Wu, C.~Feichtenhofer, T.~Darrell, and S.~Xie.
A ConvNet for the 2020s.
In \textit{Proceedings of the IEEE/CVF Conference on Computer Vision and Pattern Recognition (CVPR)}, 2022.

\bibitem{chen2025itfa}
Y.~Chen and Y.~Peng.
ITFA-DNN: Interactive time-frequency attention for environmental sound classification.
\emph{Pattern Analysis and Applications}, 2025.

\bibitem{chang2026rectangular}
L.~Chang, H.~Wang, and Z.~Li.
Rectangular kernels with self-distillation for environmental sound classification.
\emph{Pattern Analysis and Applications}, 2026.

\bibitem{crnncoordatt2026}
B.~Kim and J.~Park.
Lightweight CRNN with coordinate attention for ESC-50.
\emph{arXiv preprint}, 2026.

\bibitem{mobilenetv2spa2025}
A.~Kumar, S.~Singh, and R.~Gupta.
MobileNetV2 with spectral pooling attention and metric learning for environmental sound classification.
\emph{Expert Systems with Applications}, 2025.

\bibitem{gong2021psla}
Y.~Gong, Y.-A.~Chung, and J.~Glass.
PSLA: Improving audio tagging with pretraining, sampling, labeling, and aggregation.
In \textit{Proceedings of the IEEE/CVF International Conference on Computer Vision (ICCV)}, 2021.

\bibitem{he2016resnet}
K.~He, X.~Zhang, S.~Ren, and J.~Sun.
Deep residual learning for image recognition.
In \textit{Proceedings of the IEEE Conference on Computer Vision and Pattern Recognition (CVPR)}, 2016.

\bibitem{huang2016stochasticdepth}
G.~Huang, Y.~Sun, Z.~Liu, D.~Sedra, and K.~Weinberger.
Deep networks with stochastic depth.
In \textit{Proceedings of the European Conference on Computer Vision (ECCV)}, 2016.

\bibitem{kaplan2020scaling}
J.~Kaplan, S.~McCandlish, T.~Henighan, T.~B.~Brown, B.~Chess, R.~Child, S.~Gray,
A.~Radford, J.~Wu, and D.~Amodei.
Scaling laws for neural language models.
\emph{arXiv preprint arXiv:2001.08361}, 2020.

\bibitem{dosovitskiy2021vit}
A.~Dosovitskiy, L.~Beyer, A.~Kolesnikov, D.~Weissenborn, T.~Unterthiner,
M.~Dehghani, M.~Minderer, G.~Heigold, S.~Gelly, J.~Uszkoreit, and N.~Houlsby.
An image is worth 16$\times$16 words: Transformers for image recognition at scale.
In \textit{International Conference on Learning Representations (ICLR)}, 2021.

\bibitem{sandler2018mobilenetv2}
M.~Sandler, A.~Howard, M.~Zhu, A.~Zhmoginov, and L.-C.~Chen.
MobileNetV2: Inverted residuals and linear bottlenecks.
In \textit{Proceedings of the IEEE Conference on Computer Vision and Pattern Recognition (CVPR)}, 2018.

\bibitem{vaswani2017attention}
A.~Vaswani, N.~Shazeer, N.~Parmar, J.~Uszkoreit, L.~Jones, A.~N.~Gomez,
{\L}.~Kaiser, and I.~Polosukhin.
Attention is all you need.
In \textit{Advances in Neural Information Processing Systems (NeurIPS)}, 2017.

\end{thebibliography}
}

% ═══════════════════════════════════════════════════════════════════════════════
% Appendix
\appendix
\section{Training Recipe (Recommended)}

Table~\ref{tab:training} lists the recommended training hyperparameters for
SpectroConvNeXt on ESC-50, adapted from the ConvNeXt~V2
recipe~\cite{woo2023convnextv2}. These have not been experimentally validated.

\begin{table}[h]
  \centering
  \caption{Recommended training hyperparameters for SpectroConvNeXt on ESC-50.
  Not validated --- requires experimental confirmation.}
  \label{tab:training}
  \small
  \begin{tabular}{ll}
  \toprule
  \textbf{Hyperparameter} & \textbf{Value} \\
  \midrule
  Optimiser & AdamW ($\beta_1=0.9, \beta_2=0.999$) \\
  Learning rate & $1\times10^{-3}$ (cosine decay to $1\times10^{-5}$) \\
  Batch size & 128 \\
  Epochs & 300 (with early stopping) \\
  Warmup & 10 epochs (linear) \\
  Weight decay & 0.05 \\
  Label smoothing & 0.1 \\
  DropPath rate & 0.05 (Atto-S) to 0.30 (Tiny-S) \\
  Mixup $\alpha$ & 0.2 \\
  SpecAugment & Freq mask~$F=8$, Time mask~$T=8$, 2 masks each \\
  Random crop & Centre crop to 128$\times$80, resize to 128$\times$86 \\
  EMA decay & 0.9998 \\
  Loss & Cross-entropy + label smoothing \\
  \bottomrule
  \end{tabular}
\end{table}

\section{Data Preprocessing Pipeline}

The recommended preprocessing pipeline for ESC-50:

\begin{align*}
\text{Raw audio (44.1 kHz)} &\rightarrow \text{Resample to 22.05 kHz} \\
&\rightarrow \text{Trim/pad to 2.0 s (44100 samples)} \\
&\rightarrow \text{STFT: n\_fft}=1024, \text{hop\_length}=512 \\
&\rightarrow \text{Mel filterbank: 128 bands} \\
&\rightarrow \text{Log compression: } \log(\text{mel} + 10^{-6}) \\
&\rightarrow \text{Per-sample normalisation: } (x - \mu) / \sigma \\
&\rightarrow \text{Shape: } (1, 128, 86) \text{ (channels, freq, time)}
\end{align*}

% ═══════════════════════════════════════════════════════════════════════════════
% NeurIPS Paper Checklist
\newpage
\section*{NeurIPS Paper Checklist}

\paragraph{1. Claims.}
Do the main claims made in the abstract and introduction accurately reflect the
paper's contributions and scope?
\textbf{Yes.} The paper claims to present a reference architecture and
validation harness, which is what the artifacts deliver. All accuracy-related
claims are explicitly marked as hypotheses (unverified). No empirical
superiority is claimed.

\paragraph{2. Experiments.}
Does the paper provide sufficient detail to reproduce the experiments?
\textbf{Yes.} The validation harness is fully specified: 11 unit-test categories,
10 synthetic benchmarks, and 6 ablations, all executable via provided scripts.
Parameter counts and test statuses are reported. The training recipe is
specified but not validated (marked accordingly).

\paragraph{3. Theory.}
Does the paper make theoretical contributions?
\textbf{No.} The contribution is an architecture design and validation framework,
not a theoretical result. Design decisions are justified by empirical evidence
from prior work and structural arguments (aspect ratio matching, frequency-axis
information density).

\paragraph{4. Datasets.}
Does the paper require dataset use?
\textbf{No.} The validation harness uses synthetic (random) inputs only. The
architecture targets the ESC-50 dataset, but no ESC-50 data is loaded or
processed in the reported experiments.

\paragraph{5. Code.}
Does the paper release code?
\textbf{Yes.} The complete reference implementation is released as part of the
paper artifacts, including \texttt{model.py}, \texttt{layers.py},
\texttt{backbone.py}, \texttt{head.py}, and \texttt{frontend.py} for the
model, and \texttt{test\_model.py}, \texttt{smoke\_test.py},
\texttt{profile\_model.py}, and \texttt{run\_ablations.py} for the validation
harness.

\paragraph{6. Broader Impacts.}
Does the paper discuss potential broader impacts?
\textbf{Limited.} Environmental sound classification has applications in
surveillance, bioacoustics, and assistive technology. The reference architecture
is a general-purpose design artifact; specific application-level ethical
considerations depend on deployment context and are outside the scope of this
work.

\paragraph{7. Experimental Reproducibility.}
Are the experiments reproducible?
\textbf{Yes.} All unit tests and benchmarks are executable with the provided
code. The ablation runner verifies parameter-count consistency. Training
results are not available; the recommended training recipe is provided for
future reproduction.

\paragraph{8. Limitations.}
Are limitations explicitly discussed?
\textbf{Yes.} Section~6.3 discusses six specific limitations: no trained-model
evaluation, no baseline comparisons, synthetic-only benchmarks, heuristic FLOP
estimates, overfitting risk for large variants, and GRN without FCMAE
pretraining.

\end{document}
